\documentclass[11pt]{article}
\usepackage[a4paper,margin=25mm]{geometry}
\usepackage{amsmath,amssymb,amsthm,mathtools,booktabs,hyperref,xurl}
\setlength{\emergencystretch}{2em}
\newtheorem{theorem}{Theorem}
\newtheorem{lemma}[theorem]{Lemma}
\newtheorem{proposition}[theorem]{Proposition}
\newcommand{\Z}{\mathbb Z}\newcommand{\Q}{\mathbb Q}
\title{Rational Pell solutions with an integral shifted trace}
\author{\(\mathbb Z\mathbb Q\) Diophantine Project}
\date{Reconstructed checkpoint, 24 November 2026}
\begin{document}
\maketitle
\begin{abstract}
We classify rational solutions of a Pell norm for which
\(d-(a-2)c\) is integral. For positive \(a\ge2\) there are exactly
two unit orbits, coincident at \(a=2\); the nonintegral orbit has a
single explicit denominator. For negative parameters a finite seed
classification replaces the two-orbit statement. A reflection preserves
the exact shifted trace while interchanging the positive-parameter
integral and nonintegral branches. We also resolve the zero-parameter
edge case used in the source's nonvanishing argument. These are newly
reconstructed proofs, not a recovered copy of the previous conversation.
No integrality predicate is deleted and no numerical undecidability
bound is asserted.
\end{abstract}

\section{Denominators and the unit action}

Throughout this section let \(a\in\Z\), \(|a|\ge2\), and put
\[
 A=|a|,\qquad \Delta=a^2-1,\qquad \kappa=a-2,\qquad D=4a-5.
\]
The number \(\Delta\) is positive and nonsquare. Consider rational
\(c,d\) satisfying
\begin{equation}\label{system}
 d^2-\Delta c^2=1,\qquad T=d-\kappa c\in\Z.
\end{equation}
Substitution gives
\begin{equation}\label{quadratic}
 Dc^2-2\kappa Tc-(T^2-1)=0,
 \qquad (Dc-\kappa T)^2=\Delta T^2-D.
\end{equation}
The second expression is an integer square whenever its square root
is rational. Consequently
\begin{equation}\label{v}
 v:=Dc-\kappa T=\Delta c-\kappa d\in\Z,\qquad Dc\in\Z.
\end{equation}
In particular the reduced denominator of \(c\) divides \(|D|\)
and is odd.

Let \(\epsilon=A+\sqrt\Delta>1\). Multiplication of
\(d+c\sqrt\Delta\) by \(\epsilon\) or \(\epsilon^{-1}\)
acts, respectively, by
\begin{align}
 (d,c)&\longmapsto(Ad+\Delta c,\ d+Ac),\label{forward}\\
 (d,c)&\longmapsto(Ad-\Delta c,\ -d+Ac).\label{backward}
\end{align}
These transformations preserve the norm and are mutual inverses.
Their new shifted traces are
\begin{align}
 T_+&=(A-a+2)T+Dc,\label{traceplus}\\
 T_-&=(A+a-2)T-Dc.\label{traceminus}
\end{align}
Thus both remain integral by \eqref{v}. Their second coordinates
are
\[
 c_+=T+(A+a-2)c,\qquad c_-=-T+(A-a+2)c.
\]
For \(a\ge2\) the two multipliers are \(2a-2\) and \(2\),
both coprime to \(D\), since \(D=2(2a-2)-1\) is odd.
For \(a=-A\le-2\) they are \(-2\) and \(2A+2\), both
coprime to \(4A+5=2(2A+2)+1\). Hence each unit action preserves
the exact reduced denominator of \(c\).

There is a canonical fundamental interval for this action. Since
\(d^2=1+\Delta c^2\), the coordinate \(d\) is nonzero.
After replacing \((d,c)\) by its negative if necessary, assume
\(d>0\). Then \(\lambda=d+c\sqrt\Delta>0\). There is a
unique integer \(n\) such that
\[
                  1\le\lambda\epsilon^{-n}<\epsilon.
\]
For the corresponding pair \((d_0,c_0)\), this is equivalent to
\begin{equation}\label{interval}
       0\le c_0<1,\qquad d_0=\sqrt{1+\Delta c_0^2}>0.
\end{equation}
Indeed \(\lambda\mapsto(\lambda-\lambda^{-1})/(2\sqrt\Delta)\)
is strictly increasing, with values zero at one and one at
\(\epsilon\). The shifted trace at this representative is integral.

\section{The positive-parameter classification}

\begin{theorem}\label{positive}
For \(a\ge2\), put
\[
 \rho=\frac{2a^2-4a+3+2(a-2)\sqrt\Delta}{D}.
\]
All rational solutions of \eqref{system} are exactly
\begin{equation}\label{two}
 d+c\sqrt\Delta=\sigma\epsilon^n
 \quad\text{or}\quad
 d+c\sqrt\Delta=\sigma\rho\epsilon^n,
 \qquad \sigma\in\{1,-1\},\quad n\in\Z.
\end{equation}
At \(a=2\), \(\rho=1\), so there is one orbit and every solution
is integral. At \(a\ge3\), the two orbits are distinct. The first
has integral \(c,d\), and every second-orbit solution has
\begin{equation}\label{denominator}
                 \operatorname{den}(c)=\frac{D}{\gcd(D,3)}>1.
\end{equation}
\end{theorem}
\begin{proof}
Reduce to \eqref{interval}. Its shifted trace is
\[
 f(c_0)=\sqrt{1+\Delta c_0^2}-(a-2)c_0.
\]
It is positive, since \(\sqrt\Delta>a-2\). Moreover \(f\) is
convex on \([0,1]\), with \(f(0)=1\), \(f(1)=2\), so
\(f(c_0)\le1+c_0<2\). Integrality forces \(f(c_0)=1\).
Equation \eqref{quadratic}, with \(T=1\), therefore gives
\[
                 c_0=0\quad\text{or}\quad c_0=\frac{2(a-2)}D.
\]
The first has \(d_0=1\). The second has
\[
 d_0=1+(a-2)c_0=\frac{2a^2-4a+3}D.
\]
Both are norm-one points with integral shifted trace one. For
\(a\ge3\), the second satisfies \(0<c_0<1\), so its orbit is
distinct from the first by uniqueness of the fundamental representative.
For \(a=2\) the two points coincide. Conversely every signed unit
translate of these seeds satisfies \eqref{system}, by
\eqref{traceplus}--\eqref{traceminus}.

The first orbit is integral because the unit and its inverse have
integer coordinates. The second seed has reduced denominator
\[
 \frac{D}{\gcd(D,2(a-2))}=\frac{D}{\gcd(D,3)}.
\]
For \(a\ge3\) this exceeds one. The denominator invariance already
proved extends the formula to every signed translate.
\end{proof}

The exceptional factor three is exact, not a heuristic local bound.
For instance, \(a=5\) gives denominator five, while \(a=8\) gives
denominator nine. These include the rational examples in the existing
v102 checkpoint~\cite{Obstruction}.

\section{The negative-parameter classification}

The source currently permits a signed \(a\) before its Pell argument,
so Theorem~\ref{positive} alone is insufficient for a predicate deletion.

\begin{theorem}\label{negative}
Suppose \(a=-A\le-2\), and put \(K=4A+5\). For each integer
\(t\in\{1,\ldots,2A+1\}\) such that
\begin{equation}\label{seed-square}
                       v^2=(A^2-1)t^2+K
\end{equation}
is an integer square, take the positive square root \(v\) and set
\begin{equation}\label{negative-seeds}
 c_t=\frac{(A+2)t-v}{K},\qquad
 d_t=\frac{(A+2)v-(A^2-1)t}{K}.
\end{equation}
These are distinct fundamental-interval representatives. Every solution
of \eqref{system} has a unique form
\[
       d+c\sqrt\Delta=\sigma(d_t+c_t\sqrt\Delta)\epsilon^n,
 \qquad \sigma\in\{1,-1\},\ n\in\Z.
\]
The exact denominator on this orbit is
\[
 \operatorname{den}(c)=
 \frac{K}{\gcd\bigl(K,(A+2)t-v\bigr)}.
\]
\end{theorem}
\begin{proof}
At \eqref{interval}, the shifted trace is
\[
 g(c_0)=\sqrt{1+\Delta c_0^2}+(A+2)c_0.
\]
This function increases strictly from \(g(0)=1\) to
\(g(1)=2A+2\). Hence its integer value belongs to the stated
finite range. Equation \eqref{v} now reads
\(v=\Delta c_0+(A+2)d_0>0\) and gives \eqref{seed-square}.
Solving for \(c_0,d_0\) gives \eqref{negative-seeds}.

Conversely, fix an admissible \(t,v\). For \(t=1\), one has
\(v=A+2\) and obtains \((c_t,d_t)=(0,1)\). For \(t>1\),
\((A+2)^2t^2-v^2=K(t^2-1)>0\), so \(c_t>0\).
The formula for \(d_t\) is positive, since
\((A+2)v>\sqrt\Delta\,v>\Delta t\).
Direct substitution gives the norm one and shifted trace \(t\).
Consequently \(g(c_t)=t<2A+2=g(1)\), so \(c_t<1\).
The strict increase of \(g\) makes distinct \(t\)'s give distinct
seeds. The fundamental-interval uniqueness gives the claimed complete
orbit decomposition. The denominator formula follows directly from
\eqref{negative-seeds} and its invariance under the unit action.
\end{proof}

There need not be only two seeds. At \(a=-5\), precisely
\(t=1,2,5\) occur, giving
\[
 (c_t,d_t)=(0,1),\quad(3/25,29/25),\quad(2/5,11/5).
\]
Thus even the possible nonintegral denominators differ between the
positive and negative parameter cases.

\section{A reflection invisible to the first linear equation}

\begin{proposition}\label{reflection}
For any \(a\in\Z\), put \(\Delta=a^2-1\), \(D=4a-5\),
and \(\kappa=a-2\). If rational \(c,d\) satisfy the norm equation
and \(T=d-\kappa c\), then
\begin{equation}\label{reflection-formula}
 c^\#=\frac{2\kappa T}{D}-c,\qquad
 d^\#=\frac{2\Delta T}{D}-d
\end{equation}
satisfy the same norm and the same value of \(T\). This is an
involution. For \(a\ge3\), it interchanges the two orbits of
Theorem~\ref{positive}.
\end{proposition}
\begin{proof}
The two roots of the first equation in \eqref{quadratic} sum to
\(2\kappa T/D\). Hence \(c^\#\) is its other root. Since
\(\kappa^2+D=\Delta\), the stated formula for \(d^\#\) equals
\(T+\kappa c^\#\), proving both assertions. Applying the same root
exchange twice returns the original point.

For positive \(a\), observe that
\[
 \rho=\frac{\sqrt\Delta+\kappa}{\sqrt\Delta-\kappa},\qquad
 d^\#+c^\#\sqrt\Delta
       =\frac{\rho}{d+c\sqrt\Delta}.
\]
Thus \(\sigma\epsilon^n\) maps to
\(\sigma\rho\epsilon^{-n}\), and conversely. The two orbits
are distinct when \(a\ge3\).
\end{proof}

The first linear source equation has the form
\(T=bw+D\gamma\), with the other displayed inputs integral.
The reflection preserves its entire right side, not merely a congruence
class. That equation therefore cannot exclude the nonintegral branch.
The exact rational family in~\cite{Obstruction} additionally survives
the subsequent norm/CRT subsystem. It fails the full source's closeness
and growth conditions; neither it nor the classification is a full-source
counterexample.

\section{The zero-parameter edge case and the remaining task}

\begin{lemma}\label{zero}
If \(a=0\), \(c,d\in\Q\), \(d^2+c^2=1\), and
\(d+2c\in\Z\), then
\[
 (c,d)\in\{(0,\pm1),(\pm1,0),
 (4/5,-3/5),(-4/5,3/5),(3/5,4/5),(-3/5,-4/5)\}.
\]
In particular \(|c|<1/2\) implies \(c=0\).
\end{lemma}
\begin{proof}
Cauchy--Schwarz gives \(|T|=|d+2c|\le\sqrt5\), so
\(T\in\{-2,-1,0,1,2\}\). Substitution gives
\[
                   c=\frac{2T\pm\sqrt{5-T^2}}5.
\]
At \(T=0\) the square root is irrational. The four remaining
values give precisely the listed eight points.
\end{proof}

In the existing source proof~\cite{Signed}, the assumption \(s=0\)
forces \(a=Y=p=0\), \(k=\pm1\), and the strict strengthened
margin forces \(|c|<1/2\). Lemma~\ref{zero} therefore recovers
\(c=0\) without a c-integrality assumption. The same margin then
gives \(C_*=0\), since \(C_*\) is independently integral.
The existing contradiction with \(A_0-1\mid5^{60}\) applies.
This repairs that preliminary nonzero step if the c predicate is
tentatively removed. It does not prove any subsequent integral Pell
index assertion for the remaining rational orbits.

The precise unresolved task is to use the full small-Pell equation,
its closeness inequality, the reconstructed large source parameters,
and the remaining Pell relations to rule out all nonintegral seeds and
their translates. Neither the positive-parameter classification nor
its signed extension does that. All current numerical bounds and
integer predicates are retained.

\section{Verification scope and provenance}

This checkpoint was reconstructed on 24 November 2026 after the earlier
conversation was reported to have ended during publication. It is not
identified with any unlocated original artifact or numerical claim.
The preceding committed v102 obstruction is the comparison point.

The independent script \texttt{verify\_classification.py} checks
symbolic identities and exhaustively enumerates the fundamental
interval for a finite range of positive and negative parameters using
rational square tests. It compares that enumeration with the stated
seeds, checks exact orbit normalization, signed forward and inverse
translations, denominator invariance, the reflection, and the zero
edge case. The quantified proof is the one given above; finite checks
are not substituted for it. Run from the repository root:
{\small\begin{verbatim}
python -B \
  undecidability/pell_classification_reconstructed_v105/verify_classification.py
\end{verbatim}}

\begin{thebibliography}{9}
\small
\bibitem{Obstruction} Project, \emph{Checkpoint: the three Pell norms
do not by themselves certify c}, v102, 23 November 2026.
\par\nolinkurl{undecidability/c_integrality_checkpoint_v102/README.md}.
\bibitem{Signed} Project, \emph{Both first-Pell signs have the same
coding projection}, v101, 23 November 2026.
\par\nolinkurl{undecidability/pell_signed_branch_v101/signed_branch_projection.tex}.
\end{thebibliography}
\end{document}
